\chapter{บทนำ}

\section{ความเป็นมาและความสำคัญของปัญหา}

กระบวนการค้นหาและพัฒนายาใหม่ (Drug Discovery) ในปัจจุบันต้องเผชิญกับข้อจำกัดด้านระยะเวลาและต้นทุนมหาศาล การนำยาหนึ่งชนิดออกสู่ตลาดใช้เวลาเฉลี่ย 10--15 ปี และมีค่าใช้จ่ายหลายพันล้านบาท ขั้นตอนที่เป็นคอขวดสำคัญคือการคัดกรองสารเคมีจำนวนมหาศาลเพื่อหาสารที่มีศักยภาพ ซึ่งการทดสอบฤทธิ์ทางชีวภาพในห้องปฏิบัติการจริง (Wet Lab Experiment) กับโปรตีนเป้าหมายมีขั้นตอนที่ซับซ้อนและมีค่าใช้จ่ายสูง

เชื้อไวรัส SARS-CoV-2 ซึ่งเป็นสาเหตุของโรค COVID-19 มีเอนไซม์ Main Protease (\Mpro\ หรือเรียกอีกชื่อว่า \CLpro) เป็นเอนไซม์ชนิดซิสเตอีนโปรตีเอส (Cysteine Protease) ที่มีบทบาทสำคัญยิ่งในวงจรชีวิตของไวรัส เอนไซม์นี้ทำหน้าที่ตัดสายพอลิโปรตีนที่ไวรัสสร้างขึ้นให้กลายเป็นโปรตีนย่อยที่ทำงานได้ (Non-structural Proteins) ซึ่งจำเป็นต่อการเพิ่มจำนวนของไวรัส การยับยั้งการทำงานของ \Mpro\ จึงสามารถหยุดยั้งกระบวนการแบ่งตัวของไวรัสได้โดยตรง และเนื่องจากมนุษย์ไม่มีเอนไซม์ที่มีความคล้ายคลึงกับ \Mpro\ ทำให้เอนไซม์นี้เป็นเป้าหมายยาที่น่าสนใจอย่างยิ่ง เพราะมีโอกาสเกิดผลข้างเคียงต่อเซลล์มนุษย์ต่ำ

การพัฒนาวิธีการคำนวณทางคอมพิวเตอร์แบบจำลองในคอมพิวเตอร์ (In Silico) จึงเข้ามามีบทบาทสำคัญในการคัดกรองและประเมินประสิทธิภาพสารเคมีเบื้องต้นก่อนนำไปทดสอบจริง แนวทางหนึ่งที่ได้รับการยอมรับคือการศึกษาความสัมพันธ์เชิงปริมาณระหว่างโครงสร้างทางเคมีกับการออกฤทธิ์ทางชีวภาพ (Quantitative Structure-Activity Relationship: QSAR) ซึ่งในทางวิทยาการข้อมูลถือเป็นโจทย์ปัญหาประเภทการถดถอย (Regression Problem) โดยมีเป้าหมายเพื่อทำนายค่าความสามารถในการยับยั้งเอนไซม์ ซึ่งวัดจากค่าความเข้มข้นที่ยับยั้งการทำงานได้ครึ่งหนึ่ง (Half Maximal Inhibitory Concentration: \ICfifty) ในหน่วยนาโนโมลาร์ (nM) และแปลงให้อยู่ในรูปตัวแปรต่อเนื่องที่มีการกระจายตัวแบบสมมาตร (\pICfifty)

การทำนายฤทธิ์สารเคมีอย่างมีประสิทธิภาพจำเป็นต้องอาศัยการสกัดคุณลักษณะ (Feature Extraction) จากโครงสร้างโมเลกุลในรูปแบบ SMILES ประกอบด้วยคุณลักษณะเชิงโครงสร้างย่อย 2 มิติ (Morgan Fingerprints หรือ ECFP4) และคุณสมบัติทางกายภาพเคมีตามกฎของลิพินสกี (Lipinski's Descriptors) ได้แก่ น้ำหนักโมเลกุล (MW) ค่าสัมประสิทธิ์การละลาย (LogP) จำนวนอะตอมผู้ให้ไฮโดรเจน (HBD) และจำนวนอะตอมผู้รับไฮโดรเจน (HBA) มาประกอบเป็นชุดคุณลักษณะแบบผสม (Hybrid Features)

นอกจากนี้ ปัญหาสำคัญที่พบบ่อยในงาน QSAR คือการรั่วไหลของข้อมูล (Data Leakage) ระหว่างการปรับสเกลข้อมูลกับกระบวนการตรวจสอบความถูกต้องข้ามกลุ่ม (Cross-Validation) ซึ่งทำให้ผลลัพธ์ที่รายงานสูงเกินจริง โครงงานนี้จึงประยุกต์ใช้ Scikit-learn Pipeline ในการผนวกขั้นตอนการปรับมาตรฐานข้อมูลเข้ากับแบบจำลอง และใช้เทคนิคปัญญาประดิษฐ์ที่อธิบายได้ (Explainable AI: SHAP) เพื่อวิเคราะห์ความสำคัญของคุณลักษณะโครงสร้าง พร้อมทำการจับคู่ย้อนกลับ (Bit Mapping) เป็นภาพโครงสร้างย่อย 2 มิติ เพื่ออธิบายความสัมพันธ์ระหว่างโครงสร้างโมเลกุลกับฤทธิ์การยับยั้งเอนไซม์ได้อย่างโปร่งใสและตรวจสอบได้

\section{วัตถุประสงค์ของโครงงาน}

\begin{enumerate}[leftmargin=2em, itemsep=0.5em]
  \item เพื่อพัฒนาและเปรียบเทียบประสิทธิภาพของแบบจำลองการเรียนรู้ของเครื่อง (Machine Learning) ในการทำนายค่าฤทธิ์การยับยั้งเอนไซม์ SARS-CoV-2 Main Protease (\pICfifty) แบบ In Silico
  \item เพื่อสร้างกระบวนการเตรียมและสกัดคุณลักษณะโมเลกุลแบบผสม (Hybrid Features) จาก Morgan Fingerprints และ Lipinski Descriptors ผ่าน Scikit-learn Pipeline ที่ป้องกันการรั่วไหลของข้อมูล (Data Leakage)
  \item เพื่อวิเคราะห์และอธิบายความสัมพันธ์ระหว่างคุณลักษณะทางโครงสร้างกับฤทธิ์ทางชีวภาพ (SAR Analysis) โดยใช้เทคนิค Explainable AI (SHAP) ร่วมกับการสกัดภาพโครงสร้างโมเลกุลย่อย 2 มิติ
\end{enumerate}

\section{ขอบเขตของโครงงาน}

\begin{enumerate}[leftmargin=2em, itemsep=0.4em]
  \item มุ่งเน้นศึกษาเฉพาะสารยับยั้งเอนไซม์ Main Protease (\Mpro) ของเชื้อไวรัส SARS-CoV-2 โดยใช้ข้อมูลจากฐานข้อมูล ChEMBL
  \item ใช้ค่า \ICfifty\ ที่มีหน่วยการวัดเป็นนาโนโมลาร์ (nM) เท่านั้น และแปลงเป็น \pICfifty\ สำหรับการสร้างแบบจำลองประเภทการถดถอย
  \item สกัดคุณลักษณะโมเลกุลแบบ 2 มิติ ได้แก่ Morgan Fingerprints ขนาด 1{,}024 บิต และ Lipinski Descriptors 4 ค่า รวมเป็น 1{,}028 คุณลักษณะ
  \item เปรียบเทียบแบบจำลองการเรียนรู้ของเครื่อง 4 ชนิด ได้แก่ Dummy Regressor (Baseline), Random Forest, XGBoost และ Support Vector Regressor (SVR)
  \item ประเมินประสิทธิภาพด้วยตัวชี้วัด \Rsq, RMSE และ MAE ร่วมกับการทำ 5-Fold Cross Validation
\end{enumerate}

\section{ข้อจำกัดของโครงงาน}

\begin{enumerate}[leftmargin=2em, itemsep=0.4em]
  \item การศึกษานี้เป็นการจำลองและวิเคราะห์ทางคอมพิวเตอร์ (In Silico) เท่านั้น ไม่ได้มีการทดสอบยืนยันในห้องปฏิบัติการจริง (In Vitro Validation)
  \item การวิเคราะห์ SAR ผ่าน SHAP เป็นการระบุ \textbf{ความสัมพันธ์เชิงสถิติ (Statistical Association)} ระหว่างชิ้นส่วนโมเลกุลกับการทำนายของแบบจำลอง ไม่สามารถอ้างเป็นกลไกการจับทางชีวภาพเชิงสาเหตุระดับอะตอม 3 มิติ (Causal Mechanism) ได้
  \item ไม่ครอบคลุมการคำนวณพลศาสตร์โมเลกุล 3 มิติ (Molecular Dynamics) หรือการจำลองการจับตัวของลิแกนด์กับโปรตีน (Molecular Docking)
\end{enumerate}

\section{ประโยชน์ที่คาดว่าจะได้รับ}

\begin{enumerate}[leftmargin=2em, itemsep=0.4em]
  \item ได้แบบจำลองการเรียนรู้ของเครื่องที่มีความแม่นยำและปราศจากปัญหา Data Leakage สำหรับทำนายค่าฤทธิ์การยับยั้งเอนไซม์ SARS-CoV-2 Main Protease (\pICfifty) เบื้องต้นได้อย่างรวดเร็ว
  \item ทราบถึงคุณลักษณะทางโครงสร้างย่อย 2 มิติ (Substructures) และคุณสมบัติทางเคมีกายภาพที่มีความสัมพันธ์ต่อระดับการออกฤทธิ์ยับยั้งเอนไซม์ ซึ่งสามารถใช้เป็นแนวทางชี้นำสำหรับนักเคมียาในการออกแบบโมเลกุลตัวยาใหม่
  \item ได้ชุดข้อมูลโครงสร้างโมเลกุลมาตรฐานที่ผ่านการทำความสะอาดและสกัดคุณลักษณะแล้ว (\texttt{clean\_bioactivity\_master.csv}) เพื่อนำไปต่อยอดในงานวิจัยด้านเคมีสารสนเทศ (Cheminformatics) และการพัฒนาเว็บแอปพลิเคชันคัดกรองตัวยาในอนาคต
\end{enumerate}

\clearpage
\section{นิยามศัพท์เฉพาะ}

\begin{longtable}{L{4.2cm} L{9.6cm}}
\caption{นิยามศัพท์เฉพาะที่ใช้ในรายงาน}\label{tab:glossary}\\
\toprule
\textbf{ศัพท์} & \textbf{ความหมาย}\\
\midrule
\endfirsthead
\multicolumn{2}{l}{\small\textit{(ตารางที่ \thetable\ ต่อจากหน้าที่แล้ว)}}\\
\toprule
\textbf{ศัพท์} & \textbf{ความหมาย}\\
\midrule
\endhead
\bottomrule
\endfoot

\Mpro\ / \CLpro & Main Protease หรือ 3C-like Protease เอนไซม์ซิสเตอีนโปรตีเอสของ SARS-CoV-2 ที่ทำหน้าที่ตัดสายพอลิโปรตีนของไวรัส\\
\ICfifty & ความเข้มข้นของสารที่ยับยั้งการทำงานของเอนไซม์ได้ร้อยละ 50 ค่ายิ่งต่ำแสดงว่าสารยิ่งมีฤทธิ์แรง\\
\pICfifty & ค่าลอการิทึมติดลบของ \ICfifty\ ในหน่วยโมลาร์ คำนวณจาก $9-\log_{10}(\ICfifty[\mathrm{nM}])$ ค่ายิ่งสูงแสดงว่าสารยิ่งมีฤทธิ์แรง\\
SMILES & Simplified Molecular Input Line Entry System ระบบแทนโครงสร้างทางเคมีด้วยสายอักขระ\\
Morgan Fingerprints (ECFP4) & เวกเตอร์บิตที่เข้ารหัสโครงสร้างย่อย 2 มิติของโมเลกุล โดยพิจารณาสภาพแวดล้อมรอบอะตอมภายในรัศมีที่กำหนด\\
Lipinski Descriptors & คุณสมบัติทางเคมีกายภาพ 4 ค่า (MW, LogP, HBD, HBA) ที่บ่งบอกความเป็นไปได้ในการพัฒนาเป็นยารับประทาน\\
QSAR & Quantitative Structure-Activity Relationship การศึกษาความสัมพันธ์เชิงปริมาณระหว่างโครงสร้างเคมีกับการออกฤทธิ์\\
Data Leakage & การรั่วไหลของสารสนเทศจากชุดข้อมูลทดสอบเข้าสู่กระบวนการฝึกแบบจำลอง ทำให้ผลประเมินสูงเกินจริง\\
SHAP & SHapley Additive exPlanations วิธีอธิบายการตัดสินใจของแบบจำลองตามทฤษฎีเกม\\
Scaffold Split & การแบ่งข้อมูลโดยจัดโมเลกุลที่มีโครงสร้างแกนกลางเดียวกันไว้ฝั่งเดียวกัน เพื่อทดสอบความสามารถทำนายสารแกนโครงสร้างใหม่\\
Applicability Domain & ขอบเขตปริภูมิเคมีที่แบบจำลองมีความน่าเชื่อถือในการทำนาย\\
Activity Cliff & ปรากฏการณ์ที่โมเลกุลมีโครงสร้างคล้ายกันมากแต่มีฤทธิ์ทางชีวภาพแตกต่างกันอย่างมาก\\
\end{longtable}
